% !TEX TS-program = pdflatex
\documentclass[10pt,a4paper]{altacv}

% Marges équilibrées pour 1 page complète
\geometry{left=1.2cm,right=1.2cm,top=1.0cm,bottom=0.8cm,footskip=0pt}

\usepackage[T1]{fontenc}
\usepackage[utf8]{inputenc}
\usepackage[default]{lato}
\usepackage{microtype}
\DisableLigatures{encoding=*, family=*} % ATS-fix: désactive les ligatures en pdfLaTeX

\definecolor{Navy}{HTML}{0A2540}
\definecolor{SlateGrey}{HTML}{2E2E2E}
\definecolor{LightGrey}{HTML}{555555}
\colorlet{heading}{Navy}
\colorlet{accent}{Navy}
\colorlet{emphasis}{SlateGrey}
\colorlet{body}{LightGrey}

\renewcommand{\itemmarker}{{\small\textbullet}}
\renewcommand{\ratingmarker}{\faCircle}

\usepackage[hidelinks]{hyperref}
\hypersetup{
  pdftitle={CV -- Achraf Izelmadene -- Candidature RTE},
  pdfauthor={Achraf Izelmadene},
  pdfsubject={Stage PFE Pilotage Portefeuille Postes Électriques HTB},
  pdfkeywords={RTE, postes électriques, HTB, HTA, BT, réseaux électriques, gestion de projet, smart grid, Python}
}

% Photo avec diamètre réduit à 1.9cm et cadrage naturel
\photo{1.9cm}{assets/imageee.png}

% Style compact pour les sections
\renewcommand{\cvsection}[2][]{%
  \vspace{4pt}%
  {\color{heading}\Large\bfseries #2 \hspace{0.2em}\leaders\vrule height 3pt depth -1.5pt\hfill\mbox{}}\par\vspace{3pt}
}

\renewcommand{\cvevent}[4]{%
  {\normalsize\bfseries\color{accent}#1\par}\vspace{1pt}
  \ifstrequal{#2}{}{}{%
    {\small\itshape\color{emphasis}#2\ifstrequal{#4}{}{}{ - #4}}%
  }%
  ~\hfill
  \ifstrequal{#3}{}{}{{\small\color{emphasis}#3}}%
  \par\vspace{2pt}
}

% Header fusionné avec À Propos
\makeatletter
\newcommand{\target}[1]{\def\@target{#1}}
\newcommand{\apropos}[1]{\def\@apropos{#1}}

\renewcommand{\makecvheader}{%
  \begingroup
    \ifdef{\@photodiameter}{%
    \begin{minipage}[c]{\@photodiameter}
    \if@normalphoto
      \includegraphics[width=\linewidth]{\@photo}
    \else
      \begin{tikzpicture}
        \clip (0,0) circle (0.5\linewidth);
        \node at (0,-0.04\linewidth) {\includegraphics[width=1.15\linewidth]{\@photo}};
      \end{tikzpicture}
    \fi%
    \end{minipage}\hspace{1.2em}%
    \begin{minipage}[c]{\dimexpr\linewidth-\@photodiameter-1.5em}}{}%
    \raggedright\color{emphasis}%
    {\LARGE\bfseries\MakeUppercase{\@name}\par}\vspace{1.5pt}
    {\normalsize\bfseries\color{accent}\@tagline\par}\vspace{3pt}
    \ifdef{\@target}{{\small\textbf{\@target}\par}\vspace{2pt}}{}%
    \ifdef{\@apropos}{{\small\color{emphasis}\@apropos\par}\vspace{3pt}}{}%
    \vspace{2.5pt}%
    {\footnotesize\color{emphasis}\@personalinfo\par}
    \ifdef{\@photodiameter}{%
    \end{minipage}\par}{}%
  \endgroup\vspace{1pt}
}
\makeatother

\begin{document}

\name{Izelmadene Achraf}
\tagline{Élève Ingénieur en Génie Électrique et Systèmes Intelligents (GESI)}

% Profil adapté à l'offre RTE (Postes HTB / Direction Technique Lyon)
\target{À la recherche d’un stage PFE de 5 à 6 mois dès janvier-mars 2027 en pilotage de projets et postes électriques HTB (RTE).}
\apropos{Double compétence en génie électrique (réseaux HTA/BT, postes de livraison, bilans de puissance) et en gestion de projets / outils décisionnels (certifié 3DEXPERIENCE, dashboards Python, analyse financière). Je souhaite apporter ma rigueur et mes compétences techniques au pilotage du portefeuille d’investissements HTB chez RTE.}

\personalinfo{%
  \href{mailto:achraf.izelmadene@gmail.com}{achraf.izelmadene@gmail.com}%
  \hspace{0.9em}\textcolor{body!60}{|}\hspace{0.9em}%
  +212699098736%
  \hspace{0.9em}\textcolor{body!60}{|}\hspace{0.9em}%
  Marrakech, Maroc%
  \hspace{0.9em}\textcolor{body!60}{|}\hspace{0.9em}%
  Permis B%
}

\makecvheader

\cvsection{FORMATION}
\cvevent{Étudiant en cycle d’ingénieur en Génie Électrique et Systèmes Intelligents}{École Marocaine des Sciences de l’Ingénieur - Marrakech}{2025 - Présent}{}
{\footnotesize\color{emphasis}Conception des installations électriques HTA/BT, Réseaux électriques \& Sécurité, Énergies renouvelables : production \& stockage, Smart Grids, Audit \& Efficacité énergétique, Électronique de puissance, Entraînement électrique à vitesse variable, Supervision industrielle \& Digital Twin, Hygiène, Sécurité \& Environnement (HSE), Lean Management, Management \& Contrôle qualité, Gestion de projet...}\par\vspace{2.5pt}

\cvevent{Licence SMP, option Automatique}{Faculté des Sciences Semlalia - Marrakech}{2020 - 2024}{}
{\footnotesize\color{emphasis}Électricité, Électronique analogique, Électromagnétisme, Asservissements linéaires, Automatismes et capteurs, Traitement du signal...}\par

\cvsection{EXPÉRIENCE}
\cvevent{Stagiaire Bureau d’Études}{Core Engineering Solutions}{juil. - août 2026}{Berrchid, Casablanca}
\begin{itemize}[leftmargin=*, label=\textbullet, nosep, itemsep=1.5pt]
  \item Dimensionnement complet d’un poste de livraison HTA/BT et bilan de puissance pour une zone industrielle.
  \item Étude d’intégration PV en autoconsommation (PVsyst), schémas unifilaires, notes de calcul et DOE (AutoCAD, Caneco BT).
\end{itemize}
\vspace{2pt}
\cvevent{Stagiaire en Unité de Production}{Menara Prefa}{juil. 2024}{Marrakech}
\begin{itemize}[leftmargin=*, label=\textbullet, nosep, itemsep=1.5pt]
  \item Participation à la migration d’un automate vers Siemens S7 (TIA Portal) et suivi des opérations de production.
  \item Maintenance préventive et corrective sur lignes de production, suivi des indicateurs d’équipements.
\end{itemize}

\cvsection{PROJETS}
{\normalsize\bfseries\color{accent} Détection d’anomalies Smartgrid (Clustering \& Séries Temporelles)}\par\vspace{1.5pt}
\begin{itemize}[leftmargin=*, label=\textbullet, nosep, itemsep=1.5pt]
  \item Analyse et partitionnement de profils de consommation électrique sur réseaux Smartgrids (KMeans, DBSCAN, GMM).
  \item Développement d’un dashboard interactif (Python, Streamlit) pour le monitoring d’anomalies et l’aide à la décision.
\end{itemize}
\vspace{2pt}
{\normalsize\bfseries\color{accent} SmartStruct : Éclairage Public Piézoélectrique (Energy Harvesting \& IoT)}\par\vspace{1.5pt}
\begin{itemize}[leftmargin=*, label=\textbullet, nosep, itemsep=1.5pt]
  \item Conception d’un système de récupération d’énergie pour l’éclairage public via dalles piézoélectriques connectées.
  \item Prototype fonctionnel (PoC) sélectionné et présenté lors du congrès scientifique COCIA’26 (EST Casablanca).
\end{itemize}
\vspace{2pt}
{\normalsize\bfseries\color{accent} Commande d’un Gradateur via Arduino et Simulink (Électronique de puissance)}\par\vspace{1.5pt}
\begin{itemize}[leftmargin=*, label=\textbullet, nosep, itemsep=1.5pt]
  \item Modélisation sous Simulink et interfaçage Hardware-in-the-Loop (HIL) avec carte Arduino en temps réel.
  \item Contrôle et régulation de puissance par pilotage précis de l’angle d’amorçage.
\end{itemize}

\cvsection{COMPÉTENCES}
\begin{tabular}{@{}ll@{}}
\textbf{Réseaux \& Bureau d’études :} & Postes HTA/BT, Bilans de puissance, Notes de calcul, Caneco BT, AutoCAD, PVsyst \\[1.5pt]
\textbf{Gestion de Projet \& Outils :} & 3DEXPERIENCE (PLM/Gantt), Analyse financière (VAN, ROI), Pack Office / Excel, ProjectLibre \\[1.5pt]
\textbf{Data, IA \& Automatisation :} & Python (Pandas, Scikit-learn, Streamlit), Dashboards, Modbus RTU, TIA Portal \\[1.5pt]
\textbf{Compétences transversales :} & Rigueur d’analyse, Esprit de synthèse, Communication transversale, HSE \& Qualité
\end{tabular}

\cvsection{CERTIFICATIONS}
\begin{itemize}[leftmargin=*, label=\textbullet, nosep, itemsep=1.5pt]
  \item \textbf{Certified 3DEXPERIENCE Collaborative Industry Innovator}~\hfill {\color{emphasis}Dassault Systèmes}
  \item \textbf{Foundations of IoT Systems and Industrial Automation}~\hfill {\color{emphasis}LearnQuest (Coursera)}
  \item \textbf{Sensors and Sensor Circuit Design}~\hfill {\color{emphasis}University of Colorado Boulder (Coursera)}
  \item \textbf{MATLAB / Simulink Onramp}~\hfill {\color{emphasis}MathWorks}
\end{itemize}

\cvsection{ACTIVITÉS \& ENGAGEMENT}
\begin{itemize}[leftmargin=*, label=\textbullet, nosep, itemsep=1.5pt]
  \item \textbf{Hack The Intelligence (AISEC 2026) - 1\textsuperscript{er} Prix :} Développement d’un prototype Smart Agriculture IoT.
  \item \textbf{Vision AI Competition (EMSI 2026) - 2\textsuperscript{ème} Prix :} Conception d’un système de traduction en langue des signes (ASL).
  \item \textbf{Congrès COCIA’26 - Présentation Scientifique :} Conception du prototype d’éclairage piézoélectrique SmartStruct.
  \item \textbf{Engagement Associatif \& Bénévolat (2025) :} Soutien scolaire et encadrement pédagogique d'enfants orphelins.
\end{itemize}

\cvsection{LANGUES}
\textbf{Arabe :} Langue maternelle~\hfill \textbf{Français :} DELF B2~\hfill \textbf{Anglais :} Professionnel

\end{document}
